\section{Ejemplos prácticos y plantillas de instrucciones}

\subsection{Práctica del Prompt Template}
Un inference-time prompt ideal incluye context, task difficulty, example CoT, instrumentation hints y final instruction. Por ejemplo, la siguiente plantilla se usa con frecuencia en reasoning de varios pasos:

\begin{lstlisting}[language=Python, caption={Inference-time prompt skeleton}]
System: You are a high-precision reasoning assistant.

Input context: {problem description}
Difficulty: {mark as "hard"/"critical"}
Examples:
  - Example 1 CoT: ...
  - Example 2 CoT: ...
Instruction: Think step-by-step. If uncertain, enumerate two next steps and sample multiple candidates.
\end{lstlisting}

\begin{importantbox}{Técnicas para poner en práctica la plantilla de prompt}
Vincula difficulty como metadata dentro del prompt, para que la sampling policy ajuste automáticamente temperature y token limit; usa ejemplos de casos que expresen el contraste «wrong → right», para recordarle al modelo que debe extender el reasoning en profundidad.
\end{importantbox}

\subsection{Self-Consistency Orchestration}
En la etapa de sampling, primero se generan continuamente entre 20 y 40 responses con high-temperature; en la etapa de search, estas responses se toman como seed y se activa beam/tree search para extender más las promising branches; en la etapa de iteration, se ejecutan trace feedback y la puntuación del verifier, y después se elige el answer final con base en majority vote/consistency.

\begin{knowledgebox}{Orchestration checklist}
1. Sampling: registra candidate count + diversity; 2. Search: track branch depth, stuck rate; 3. Iteration: log trace iteration count, verification score, self-correction delta.
\end{knowledgebox}

\subsection{Resumen de la sección}
Los ejemplos prácticos ayudan a concretar la idea abstract de inference-time: podemos apoyarnos en el prompt skeleton, las orchestrations de las tres etapas sampling/search/iteration, y el instrumentation checklist para escribir el pipeline como un template reutilizable.

